\chapter{Implementation}

\section{Technology Stack}

\begin{table}[H]
\centering
\caption{Technology Stack Used in SENTINEL AI}
\label{tab:tech_stack}
\begin{tabular}{l l l}
\toprule
\textbf{Category} & \textbf{Technology} & \textbf{Version} \\
\midrule
Language & Python & 3.11+ \\
Backend Framework & FastAPI & $\geq$0.100 \\
ORM & SQLAlchemy & $\geq$2.0 \\
Database & SQLite & Built-in \\
ML Framework & Scikit-learn & $\geq$1.3 \\
Gradient Boosting & XGBoost, LightGBM & $\geq$2.0, $\geq$4.0 \\
Packet Capture & Scapy & $\geq$2.5 \\
Class Balancing & imbalanced-learn (SMOTE) & $\geq$0.11 \\
Explainability & SHAP & $\geq$0.42 \\
Dashboard & Streamlit & $\geq$1.27 \\
Visualization & Plotly, Matplotlib & $\geq$5.17, $\geq$3.8 \\
Containerization & Docker, Docker Compose & Latest \\
\bottomrule
\end{tabular}
\end{table}

\section{Data Preprocessing Module}

The \texttt{preprocessing/pipeline.py} module implements the \texttt{DataPipeline} class, which handles data loading, synthetic generation, cleaning, scaling, and SMOTE balancing.

\subsection{Synthetic Data Generation}

When the CICIDS2017 dataset is unavailable, the system generates a synthetic dataset of 10,000 records that faithfully replicates the statistical distributions:

\begin{lstlisting}[caption={Synthetic CICIDS2017-like Data Generation}, label={lst:synthetic}]
def _generate_synthetic_data(self):
    np.random.seed(42)
    rows = 10000
    data = {
        'flow_duration': np.random.randint(100, 1000000, rows),
        'fwd_packets_tot': np.random.randint(1, 100, rows),
        'bwd_packets_tot': np.random.randint(1, 100, rows),
        'flow_byts_s': np.random.uniform(0, 100000, rows),
        'flow_pkts_s': np.random.uniform(0, 1000, rows),
        'pkt_len_mean': np.random.uniform(40, 1500, rows),
        'pkt_len_std': np.random.uniform(0, 500, rows),
        'iat_mean': np.random.uniform(0, 10000, rows),
        'iat_std': np.random.uniform(0, 5000, rows),
        'syn_flag_cnt': np.random.randint(0, 2, rows),
        'ack_flag_cnt': np.random.randint(0, 2, rows),
        'psh_flag_cnt': np.random.randint(0, 2, rows),
        'label': np.random.choice(['BENIGN', 'BOTNET'],
                                   rows, p=[0.9, 0.1])
    }
\end{lstlisting}

\subsection{SMOTE Application}

The class imbalance (90\% BENIGN, 10\% BOTNET) is addressed using SMOTE:

\begin{lstlisting}[caption={SMOTE Class Balancing}, label={lst:smote}]
def balance_and_split(self, X, y):
    X_train, X_test, y_train, y_test = train_test_split(
        X, y, test_size=0.2, random_state=42, stratify=y)
    smote = SMOTE(random_state=42)
    X_train_res, y_train_res = smote.fit_resample(
        X_train, y_train)
    return X_train_res, X_test, y_train_res, y_test
\end{lstlisting}

\section{Model Training Module}

The \texttt{training/trainer.py} module implements the \texttt{ModelTrainer} class, which benchmarks three classifiers:

\begin{lstlisting}[caption={Multi-Model Benchmarking}, label={lst:trainer}]
class ModelTrainer:
    def __init__(self):
        self.models = {
            "RandomForest": RandomForestClassifier(
                n_estimators=100, random_state=42),
            "XGBoost": XGBClassifier(
                use_label_encoder=False,
                eval_metric='logloss', random_state=42),
            "LightGBM": LGBMClassifier(random_state=42)
        }
\end{lstlisting}

The champion model is selected based on the highest F1-Score and serialized using Joblib. In our experiments, \textbf{XGBoost} was consistently selected as the champion model.

\section{Backend API Implementation}

The FastAPI backend (\texttt{backend\_api/main.py}) provides six endpoints. The core prediction endpoint:

\begin{lstlisting}[caption={Inference Endpoint Implementation}, label={lst:predict}]
@app.post("/predict")
async def predict_flow(features: FlowFeatures):
    df = pd.DataFrame([features.dict()])
    X_scaled = scaler.transform(df)
    prediction = model.predict(X_scaled)[0]
    probabilities = model.predict_proba(X_scaled)[0]
    return {
        "is_botnet": bool(prediction),
        "confidence": float(probabilities[prediction]),
        "timestamp": datetime.now().isoformat()
    }
\end{lstlisting}

\subsection{Inference Sequence Diagram}

The following sequence diagram illustrates the step-by-step process of flow prediction:

\begin{figure}[H]
\centering
\resizebox{0.85\textwidth}{!}{%
\begin{tikzpicture}[
    node distance=2.5cm,
    actor/.style={draw=black!70, thick, rounded corners, fill=purple!5, minimum width=2.5cm, minimum height=1cm, align=center},
    process/.style={draw=black!70, thick, rounded corners, fill=blue!5, minimum width=2.5cm, minimum height=1cm, align=center},
    model/.style={draw=black!70, thick, rounded corners, fill=green!5, minimum width=2.5cm, minimum height=1cm, align=center},
    line/.style={draw, thick, -latex'}
]

% Nodes
\node[actor] (sniffer) {Scapy Sniffer};
\node[process, right=of sniffer] (api) {FastAPI \texttt{/predict}};
\node[process, right=of api] (scaler) {StandardScaler};
\node[model, right=of scaler] (xgb) {XGBoost Model};
\node[actor, below=of api, yshift=-1cm] (db) {SQLite Database};

% Edges
\draw[line] (sniffer) -- node[above, font=\footnotesize] {1. POST JSON Flow} (api);
\draw[line] (api) -- node[above, font=\footnotesize] {2. Transform} (scaler);
\draw[line] (scaler) -- node[above, font=\footnotesize] {3. Scaled Vector} (xgb);
\draw[line] (xgb.south) -- +(0,-0.8) -| node[above right, font=\footnotesize] {4. Probability \& Class} (api.south);
\draw[line] (api) -- node[left, font=\footnotesize] {5. Async Log} (db);
\draw[line] (api.south west) -- node[above left, font=\footnotesize] {6. Return JSON} (sniffer.south east);

\end{tikzpicture}
}
\caption{Real-Time Inference Sequence}
\label{fig:inference_sequence}
\end{figure}

\subsection{SHAP Explanation Endpoint}

The \texttt{/explain/\{id\}} endpoint reconstructs the feature vector from the database and computes SHAP values using \texttt{TreeExplainer}:

\begin{lstlisting}[caption={SHAP Integration in API}, label={lst:shap_api}]
explainer = shap.TreeExplainer(model)
shap_values = explainer.shap_values(X_scaled)
return {
    "base_value": float(base_val),
    "values": val,
    "feature_names": feature_order,
    "actual_features": threat.features
}
\end{lstlisting}

\section{Real-Time Sniffer Engine}

The \texttt{realtime/sniffer.py} module implements the \texttt{RealTimeEngine} class with two operational modes:

\subsection{Live Capture Mode}

Uses Scapy's \texttt{sniff()} for real-time packet capture with biflow aggregation:

\begin{lstlisting}[caption={Biflow-Based Packet Processing}, label={lst:sniffer}]
def packet_callback(self, pkt):
    if IP in pkt:
        ip_src = pkt[IP].src
        ip_dst = pkt[IP].dst
        flow_id = tuple(sorted([
            (ip_src, sport), (ip_dst, dport)
        ])) + (proto,)
        flow = self.flows[flow_id]
        flow['pkt_count'] += 1
        flow['byte_count'] += len(pkt)
        flow['packet_lengths'].append(len(pkt))
        if flow['pkt_count'] >= 10:
            self._analyze_flow(flow_id, flow)
\end{lstlisting}

\subsection{Autonomous Simulation Mode}

When administrative privileges are unavailable, the engine automatically falls back to a high-fidelity simulation that generates realistic attack patterns (Mirai, C2 Beacon, UDP Flood, SSH Brute Force).

\section{Dashboard Implementation}

The Streamlit dashboard (\texttt{frontend\_streamlit/app.py}) implements five views with Plotly visualizations. The styling uses a custom CSS injection for the ``Cosmic-Glass'' aesthetic:

\begin{lstlisting}[caption={Cosmic-Glass CSS Injection (excerpt)}, label={lst:css}]
div[data-testid="stMetric"] {
    background: rgba(255, 255, 255, 0.03);
    backdrop-filter: blur(10px);
    border: 1px solid rgba(255, 255, 255, 0.1);
    border-radius: 15px;
    transition: transform 0.3s ease;
}
\end{lstlisting}

\section{Orchestration and Deployment}

The \texttt{run\_all.py} orchestrator launches all three services (Backend, Sniffer, Dashboard) in parallel subprocesses. Docker deployment is supported via two Dockerfiles and a \texttt{docker-compose.yml} configuration.
